\chapter{Discussion and Threats to Validity\label{cha:chapter6}}

This chapter synthesizes the empirical findings, discusses practical prompt-policy implications, and analyzes threats to validity.

\section{Main Interpretation\label{sec:main-interpretation}}

The initial interpretation is balanced. Prompt framing is not a universal tool for accuracy, but it is indeed an engineering variable in action. When considering the LLM as a system, it is important to evaluate the prompt, not view it as the final word in LLM operations. The latest studies of multilingual LLMs indicate that prompt sensitivity differs depending on the language, which adds to the importance of the prompt parameters \cite{zhu2024multilingual}.

The effects of prompts are viewed more as interface effects between the model and the evaluator than as changes in model competence. Polite code prompts lead to an increase in explanatory language, making results harder to extract, while structured wrappers decrease verbosity. In local math tasks, the compactness introduced by some prompt types is associated with poorer quality of answers. The main takeaway is that prompts, parsers, validators, and task descriptions together create a measuring system. Recent research on reducing hallucinations in LLMs provides more methods of ensuring the reliability of results \cite{tonmoy2024hallucination, gao2023retrieval}.

\section{Implications for LLM-Based Coding Systems}

Reliable coding assistants need engineering beyond prompt design: explicit output contracts and automated validation. Workflows should separate reasoning from the final artifact, produce a clean code artifact at the end, and avoid relying on prompt tricks like politeness or role labeling for correctness. The debate continues about whether ChatGPT and similar tools are genuinely beneficial for software engineering or introduce new failure modes.
The API Python results show the value of structured wrappers in some coding workflows: raw pass-rate contrasts improve under JSON and XML while tokens and duration decrease, and the mixed-effects model supports both effects. This suggests that structured wrappers can improve compatibility with automated evaluation when they align with the task and model. The local math result shows the danger of turning this into universal advice.

The same principle applies to agentic workflows, where model outputs are consumed by downstream systems \cite{liu2023agentbench, wang2024agent, xi2023rise}. Recent work on tool learning with foundation models further demonstrates that LLMs must interface with external systems while maintaining output reliability \cite{qin2023tool}. A response that provides human-friendly explanations may be unusable for automated tools due to extra commentary or formatting inconsistencies, while machine-friendly output may be too terse for human debugging. Prompt policy must therefore be directed at the actual output consumer: the instruction should state the required artifact, the evaluator should accept common valid forms, and systems serving both human and automated consumers should separate these phases rather than asking one prompt to satisfy incompatible goals.

\section{Why Code and Math Differ}

Programming and mathematical evaluations impose different requirements for output results. Code evaluation favors artifacts that are executable and compatible with the parser. Extra explanations become barriers by interfering with extraction. Structured wrappers help to signify the nature of the required output and reduce irrelevant text.

Mathematical evaluation, on the other hand, requires a correct answer derived from the reasoning process. The output contract asks for just the answer; therefore, extraction is easier, although very compact framing may coincide with less successful answer derivation. In this experiment, JSON framing lowers the number of local math tokens and reduces generation duration, while the pass rate drops. This does not prove that visible reasoning is required; it only shows that format constraints interact with task types in ways that affect downstream accuracy.

\section{Reconsidering Prompt-Engineering Folklore}

The results offer a mixed perspective on prompt folklore. The principle of "being polite" as a prompt technique was found to provide no improvement in correctness; in fact, it resulted in larger token counts and increased response times without any benefit to the pass rate for local code. Thus, the introduction of the expert role as a prompt adjustment turns out to offer little value, since it imposes only about a one percentage point penalty on local code while having an adverse effect on API math. Finally, JSON provides a mixed picture, with certain advantages for API Python and certain disadvantages for local math.

The findings of this study are therefore highly situational. Prompt effects vary depending on model settings, the type of task (e.g., coding vs. math), and the criteria for evaluation (e.g., pass rate vs. behavior). The situational character of the findings also explains contradictions in the literature. Different prompt versions implement several changes simultaneously, which, combined with the fact that different models have different levels of verbosity and sensitivity to wrappers, means that each component influences the output differently. Factorial studies investigate these components but cannot show that any single component works in all cases. Therefore, it is necessary to verify prompt practices locally, describe them accurately, and evaluate them in terms of both correctness and behavior.

\section{Practical Prompt-Policy Recommendations}

The safest default is a neutral, explicit prompt stating the task and expected output, without social or pressure cues. Use role prompts only when they reflect real task constraints and are tested. Structured wrappers suit coding workflows needing compact outputs, but not necessarily mathematical reasoning. Aggressive pressure cues should not be defaults: they show no reliable correctness gains and set a poor interaction norm. Practical evaluation must track pass rate, parser success, tokens, duration, result-only compliance, empty responses, and failure modes. Optimize safely by starting with a neutral baseline, changing one factor at a time, testing on actual task distributions, logging behavior alongside correctness, inspecting parser failures, and rerunning evaluations whenever models, providers, settings, or evaluators change. This modest recommendation argues that prompt rules don't transfer: API wins may fail locally, and token savings can hurt accuracy. A local contract beats a universal recipe for real LLM failures. Engineering is empirical: propose, test against a baseline on correctness/behavior, and decide on evidence, not style.

The local-contract view creates maintenance debt: prompt policies scatter across scripts, wrappers, system prompts, notebooks, and agent configs, becoming hidden infrastructure. Model, parser, or task changes can silently invalidate them. Treat prompt templates as version-controlled system components—track validated baselines, regressions, and model/evaluator identifiers. This replaces preference-based decisions with measured trade-offs and simplifies diagnosing failures when prototypes become production agents.

\section{Threats to Validity\label{sec:threats-to-validity}}

\subsection{Internal Validity}

The crossed design maintains internal validity because every base task appears under every prompt variant. This eliminates confounding between task difficulty and prompt condition. However, prompt variants also differ in length and framing. In practice, polite, aggressive, role, and structured-wrapper prompts can vary in token count. Consequently, some effects may reflect prompt length or prompt density rather than tone or format.

Problems relating to hardware variability, API variability, and timeouts should also be kept in mind. Empty responses have to be kept in order to maintain the factorial design; however, their causes cannot be determined from the collected output. The current thesis describes the running configurations and execution dates; retry strategy and local hardware characteristics are included in the replication records.

The crossed design restricts the order of tasks and environmental confounds, but does not fully remove them. Each basic task is performed under each prompt variant to avoid comparing easy tasks in one condition with hard tasks in another. However, the local system load and the hosted API serving conditions might be different during different calls. The measured durations should therefore be treated not as an independent characteristic of the prompt, but as observed system latency in a particular execution environment. While correctness and compliance are not greatly affected by timing noise, timeouts as well as empty responses show that runtime behavior remains part of the evaluated system.

\subsection{Construct Validity}

Pass@1 and test pass rates can't ensure complete code quality. It might happen that even if the application passes the tests, it can be weak, inefficient, or badly styled. On the other hand, a correct answer can be delivered, but it can still be unsatisfactory in terms of extraction. The shortcomings mentioned are not only imperfections; they reflect the engineering realities of an automated LLM pipeline.

The accuracy of the math is determined by how the final answer is extracted and normalized. Distortion of the results occurs when extraction is too lenient or strict. The metric based on result-only compliance is yet another constructed measure and thus may have different meanings for mathematics and code. Similarly, the API Python compliance heuristic is still not identical to the local code extraction and/or compliance rule, so comparisons between the various groups in terms of their compliance will be largely descriptive.

In addition, task descriptions also impact construct validity. While coding tasks deal with functional correctness as specified by the tests, this does not guarantee that maintainability, security, style, and generalization are taken into account. Even though tests would pass, it should still be noted that hallucinations in the generated code remain an issue \cite{zhang2025hallucination}. Math tasks are concerned with the result as extracted from the final answer rather than how well or faithfully it is explained. These decisions work perfectly for a prompt-based experiment, but they shouldn’t be confused with full evaluations of programming skill or mathematical reasoning. Hence, the thesis is focused mainly on the outcomes instead of broad statements about the capabilities of the model. Furthermore, bias in model outputs—including gender and racial bias—represents another dimension of quality that pass@1 cannot capture.

\subsection{Statistical Conclusion Validity}

The experiment includes repeated observations by base task and model. The main inferential analysis therefore uses mixed-effects logistic models with a random intercept per base task, which model this dependence structure explicitly, supported by clustered linear-probability robustness checks. Simple pairwise tests and descriptive percentage-point contrasts are retained only as the interpretability layer; they do not fully model the dependence structure on their own. Furthermore, many pairwise comparisons are exploratory, and uncorrected p-values should not be overinterpreted. Model-specific contrasts require false-discovery-rate correction before they can support strong inferential claims.

There are different levels of evidence presented. Simple contrast tables are useful for descriptive direction and effect size. Local model-specific contrasts are corrected with the Benjamini-Hochberg method because of the large number of comparisons made. Mixed-effects logistic models provide the main inference for pass/fail, with clustered linear-probability models as a fixed-effect robustness check. Results such as API Python XML and local math JSON are stronger because they are directionally stable across the mixed model, the clustered check, and the dominant-model sensitivity panels; marginal findings such as API Python JSON, which loses significance when DeepSeek V4 Flash is removed, should be read as suggestive rather than definitive.

One remaining limitation is the estimation method: the mixed-effects logistic models were fitted using variational Bayes rather than full maximum-likelihood or Markov-chain Monte Carlo estimation, which might result in somewhat understated posterior uncertainty. However, because the mixed-effects results show close agreement with the clustered-LPM robustness check—which utilizes an entirely different estimation procedure—the chosen estimation method does not alter any substantive conclusions.

\subsection{External Validity}

In the given set of tasks, there are 50 coding tasks as well as 50 mathematics tasks per model group. The tasks are domain-specific, in the sense that local tasks are derived from the national exam in Vietnam while the API tasks are from MathNet competitions. Since several coding tasks originate from mathematical problems converted into Python, the study evaluates controlled initial code generation rather than repository-level software development, and the results may be different for other programming languages, larger software engineering benchmarks, real repository tasks, multi-file or iterative workflows, or other math domains. Domain-specific applications in fields such as healthcare \cite{singhal2025medical, amugongo2025rag} and medical diagnostics \cite{chow2025medicalchatbots, ansari2025ecg} demonstrate that prompt sensitivity may vary substantially across different application domains.

The set of models consists of six local Ollama variants and two hosted API models. This makes it possible to compare local settings with hosted models but is not representative of all LLMs. Local quantized models might behave differently from models deployed using full precision. API model behavior may change over time, and therefore it is important to have accurate information on identifiers and run dates to achieve reproducibility.

This research does not claim that the selected prompt factors exhaust the prompt-design space. There may be many other prompt-related factors like few-shot examples, chain-of-thought instructions, explicit schemas, constrained decoding, tool-use instructions, or multi-turn repair prompts. Recent work on self-play fine-tuning demonstrates that iterative self-improvement can convert weak models into strong ones \cite{chen2024self}. All of these factors may interact with tone, role, and wrapper format in different ways.

The external validity of the experiment is also restricted due to the one-shot evaluation setting. Real coding assistants often operate over multiple turns, inspect test failures, edit files, or call tools. Such feedback loops may reduce some prompt-framing effects since the system can repair malformed outputs \cite{chen2022codet, madaan2023selfrefine}, or they may introduce new effects because the prompt also controls how the model responds to tool feedback. Our research is best understood as evidence about initial generation under controlled prompt variation rather than as a full model of long-running software agents.

\subsection{Evaluation Validity and Audit Evidence}

The high local math pass rate for Nemotron is likely to be a concern for reviewers. The audit of the repository data integrity indicates that there is no row-level explanation for the 91.2\% result. The high pass rate can be verified by checking the row counts, task balance, prompt-variant balance, duplicate checks, independent pass recomputation, and prompt-artifact leakage checks. Having all of these parameters match the expected design and stored pass labels strengthens confidence that the result is not caused by a copied-file mismatch, duplicate rows, arithmetic error, or obvious prompt leakage.

However, the audit has an important limitation. The math runner does not save the entire raw model response but only the normalized final-answer text in \texttt{response\_text}. The completion-token counts are still based on the raw model response. The Nemotron result should be understood as a correctness check of the final answer according to the numeric evaluator instead of proof of the unambiguity or well-formedness of the original response. It is desirable that, in the next replication, raw answers be saved along with normalized final answers for independent checking of final-answer accuracy, formatting, and any visible solution work.
